\section{\label{sec:simulation}Numerical simulation and data}

We use the mean--variance portfolio objective defined in
Eq.~\eqref{eq:portfolio-cost}. The asset universe consists of the tickers
AAPL, MSFT, NVDA, KO, JPM, AMZN, GOOGL, and META, using historical daily prices
over the period from 2021-05-27 to 2026-05-21. The data are divided into four
market windows, specified below. Unless otherwise stated, the risk-aversion
parameter is fixed at \(q=0.5\). Daily prices are converted into daily
returns, and the return vector and covariance matrix are annualized according
to
\begin{equation}
    \mu = 252\,\mu_{\mathrm{daily}},
    \qquad
    \Sigma = 252\,\Sigma_{\mathrm{daily}},
\end{equation}
where 252 is the standard approximation for the number of trading days in one
year.

For each instance, the exact optimum is obtained by enumerating the feasible
set. The main diagnostics are the probability of sampling an optimal feasible
portfolio,
\begin{equation}
    P_\star = \sum_{x:C(x)=C_{\min}} |\braket{x}{\psi}|^2,
\end{equation}
the probability of sampling a portfolio within five percent of the feasible
cost range above the optimum,
\begin{equation}
    P_{\leq 5\%}
    =
    \sum_{x:\, [C(x)-C_{\min}]/[C_{\max}-C_{\min}]\leq 0.05}
    |\langle x|\psi\rangle|^2 ,
    \label{eq:p-five-percent}
\end{equation}
\gabriel{Another potentially informative idea that is more global: plotting the cumulative probability distribution of solution probabilities as a function of the energy / gap / approximation ratio. As QAOA variants often have output distributions with significant variance, this helps us have a global notion of what's happening in terms of convergence to the ideal distribution. If a metric is needed instead of a plot, we can use some sort of normalized integral of this plot.}
and the normalized cost gap
\begin{equation}
    \Delta_{\rm norm}
    =
    \frac{\langle C\rangle-C_{\min}}{C_{\max}-C_{\min}} .
    \label{eq:normalized-gap}
\end{equation}
Here \(C_{\min}\) and \(C_{\max}\) are computed within the feasible set.
Smaller values of \(\Delta_{\rm norm}\) indicate better performance. These quantities are well defined for non-degenerate feasible cost ranges,
\(C_{\max}>C_{\min}\), which is the case for all instances considered here.

We first benchmark a five-asset instance with cardinality \(k=3\). This
instance has
\begin{equation}
    |\mathcal{F}|=\binom{5}{3}=10
\end{equation}
feasible portfolios, compared with \(2^5=32\) states in the full computational
basis. The optimal bit string is \(10101\), with
\(C(x_\star)=-0.490790\). This single-instance benchmark is used primarily to
make explicit the relation between the feasible graph and the resulting
sampling distribution.

We then perform an extended benchmark over four market windows,
\begin{equation}
\begin{array}{ll}
W_1: & 2025\text{-}02\text{-}24\;\text{to}\;2026\text{-}05\text{-}21,\\
W_2: & 2023\text{-}11\text{-}21\;\text{to}\;2025\text{-}02\text{-}21,\\
W_3: & 2022\text{-}08\text{-}24\;\text{to}\;2023\text{-}11\text{-}20,\\
W_4: & 2021\text{-}05\text{-}27\;\text{to}\;2022\text{-}08\text{-}23.
\end{array}
\end{equation}
The tested feasible spaces are
\[
    (n,k)=(5,2),(5,3),(6,2),(6,3), (6,4),
    (7,3),(7,4),(8,3),(8,4),
\]
\gabriel{Personally, I find this too small to draw conclusions related to scalability. Can we try running bigger instances (say, up to 20 assets)? I can help with this.}
with circuit depths \(p=1,2,3\). For each method and depth, the variational parameters are optimized using
COBYLA as implemented in SciPy. Each optimization is repeated over
\(N_{\rm seed}=5\) random initializations, with initial angles sampled
independently from the uniform distribution on \([0,\pi]\). Unless otherwise
stated, all methods are given the same maximum number of optimizer iterations,
\(N_{\rm iter}=120\), in the extended benchmark. For the penalty-QAOA baseline,
the penalty parameter is scanned over
\begin{equation}
    \lambda\in\{0.01,0.1,1,10,100\},
\end{equation}
and the best result within this scanned range is reported. This gives the
penalty baseline the benefit of penalty tuning within the tested range.
\gabriel{Something else we could try is using qiskit's standard conversion functions that choose this automatically, e.g. QuadraticToQubo, in case none of the selected penalties give reasonable results.}

All feasible QWOA mixers are simulated directly in the feasible subspace and
therefore have \(P_{\rm feasible}=1\) by construction. By contrast, the
transverse-field penalty-QAOA baseline evolves in the full computational basis,
so the feasibility probability is reported separately when informative.